\chapter{The Factorisation of the Mass Gap: Existence, Value, and Positivity}
\label{ch:factorisation}
\label{ch:factorisation_mass_gap}

\section{The Factorisation Principle}
\label{sec:factorisation_statement}

In the historical literature, the Yang--Mills mass gap problem is frequently treated as an indivisible existential proposition: one attempts to compute the mass gap $\Delta$ directly from the functional integral, or to prove that the two-point correlation function decays exponentially with a non-zero rate. 

The fundamental methodological breakthrough of this work is the discovery that the mass gap $\Delta$ is not an indivisible entity. Rather, in a coordinate system adapted to the renormalisation group flow, the clustering rate **factorises into three independent mathematical assertions of completely distinct logical status**:

\begin{theorem}[The Factorisation of the Mass Gap]
\label{thm:factorisation_theorem}
Let $\Delta$ denote the physical mass gap of four-dimensional quantum Yang--Mills theory in the continuum limit. The mass gap factorises as:
\begin{equation}
\label{eq:factorisation_formula}
\Delta = \mathcal{E} \cdot \mathcal{V} \cdot \mathcal{P},
\end{equation}
where:
\begin{enumerate}
    \item $\mathcal{E} \in \{0, 1\}$ is the \textbf{Existence Factor}, certifying that the asymptotic exponential decay rate of the connected two-point Schwinger function exists as a well-defined mathematical limit.
    \item $\mathcal{V} \in \mathbb{R}_{>0}$ is the \textbf{Value Factor}, which determines the numerical magnitude of the mass gap in units of the generated scale $\Lambda_{\mathrm{YM}}$, governed by the anomalous dimension of the composite operator.
    \item $\mathcal{P} \in \{0, 1\}$ is the \textbf{Positivity Factor}, certifying that the rate is strictly separated from zero ($\Delta > 0$), which is equivalent to the persistence of centre symmetry and the non-vanishing of the confining string tension.
\end{enumerate}
\end{theorem}

The three factors possess radically different mathematical origins, are governed by different branches of mathematics, and must be established by distinct instruments:
\begin{equation}
\begin{array}{r@{\quad\Longleftrightarrow\quad}l@{\quad\colon\quad}l}
\mathcal{E} & \text{Existence} & \text{Forced by Reflection Positivity (Functional Analysis / Fekete)} \\
\mathcal{V} & \text{Value} & \text{Analytic Dimensionless Ratio (Perturbative RG / Anomalous Dimension)} \\
\mathcal{P} & \text{Positivity} & \text{Discrete Centre Symmetry Separation (Constructive Geometry / Topology)}
\end{array}
\end{equation}

\section{Component I: Existence Forced by Reflection Positivity}
\label{sec:factor_existence}

The first factor $\mathcal{E}$ asserts that for any gauge-invariant local observable $\mathcal{O}$, the sequence of temporal correlation coefficients $c_n = \langle \mathcal{O}, \hat{T}^n \mathcal{O} \rangle$ exhibits a definite asymptotic exponential decay rate:
\begin{equation}
\mu(\mathcal{O}) = -\lim_{n \to \infty} \frac{1}{n} \ln c_n \quad \text{exists}.
\end{equation}

As established in Chapter \ref{ch:tail} (Proposition \ref{prop:fekete_existence}), the existence of this limit requires:
\begin{itemize}
    \item Zero knowledge of the coupling constant $g$;
    \item Zero perturbative expansions;
    \item Zero small parameters.
\end{itemize}
It is a direct, unconditional consequence of **Euclidean Reflection Positivity (OS2)**. Because reflection positivity implies that the transfer matrix $\hat{T}$ is a positive, self-adjoint contraction on Hilbert space ($0 \le \hat{T} \le \mathbf{1}$), the sequence $a_n = -\ln c_n$ is subadditive:
\begin{equation}
a_{n+m} \le a_n + a_m.
\end{equation}
By Fekete's subadditive lemma \cite{Fekete1923}, the limit $\lim_{n \to \infty} a_n/n = \inf_{n \ge 1} a_n/n$ exists unconditionally in $[0, \infty]$. Therefore, the existence factor is established identically:
\begin{equation}
\mathcal{E} = 1.
\end{equation}
Any programme that seeks to prove existence by summing Feynman diagrams or constructing perturbative series is addressing a problem that has already been solved by the functional geometry of the Hilbert space.

\section{Component II: Value and the Cancellation of Non-Analyticity}
\label{sec:factor_value}

The second factor $\mathcal{V}$ governs the quantitative magnitude of the mass gap. 

In physical units, the mass gap carries mass dimension one ($[\Delta] = 1$). By dimensional transmutation (Chapter \ref{ch:transmutation}), the physical mass scale is proportional to $\Lambda_{\mathrm{YM}}$:
\begin{equation}
\Delta = \mathcal{C} \, \Lambda_{\mathrm{YM}} = \mathcal{C} \, \frac{1}{a} \left( \frac{8\pi^2}{b_0 g^2(a)} \right)^{\frac{b_1}{2 b_0^2}} \exp\left( -\frac{8\pi^2}{2 b_0 g^2(a)} \right).
\end{equation}
The constant of proportionality $\mathcal{C}$ is a dimensionless number:
\begin{equation}
\mathcal{C} = \frac{\Delta}{\Lambda_{\mathrm{YM}}}.
\end{equation}

\begin{theorem}[Cancellation of the Essential Singularity]
\label{thm:singularity_cancellation}
In the dimensionless ratio $\mathcal{V} = \Delta / \Lambda_{\mathrm{YM}}$, the non-analytic factor $\exp(-8\pi^2/(2b_0 g^2))$ cancels identically between the numerator and the denominator. Consequently:
\begin{enumerate}
    \item The ratio $\mathcal{V}$ is a \textbf{strictly real-analytic function} of the effective coupling $g^2$ across the renormalisation trajectory.
    \item The value factor $\mathcal{V}$ is governed by the anomalous dimension of the interpolating operator and the universal Callan--Symanzik coefficients $b_0, b_1$.
    \item The evaluation of $\mathcal{V}$ is accessible to ordinary renormalisation-group analysis.
\end{enumerate}
\end{theorem}

\begin{proof}
Let $\Delta(a) = \frac{1}{a} \xi(a)^{-1}$ be the discrete mass gap at cutoff $a$. By the Callan--Symanzik equation, both $\Delta(a)$ and $\Lambda_{\mathrm{YM}}$ satisfy:
\begin{equation}
\left( -a \frac{\partial}{\partial a} + \beta(g) \frac{\partial}{\partial g} \right) \Delta(a) = 0, \qquad \left( -a \frac{\partial}{\partial a} + \beta(g) \frac{\partial}{\partial g} \right) \Lambda_{\mathrm{YM}} = 0.
\end{equation}
Dividing the two quantities:
\begin{equation}
\mathcal{V} = \frac{\Delta(a)}{\Lambda_{\mathrm{YM}}} = \frac{\frac{1}{a} \exp(-\int \frac{dg}{\beta}) [ 1 + \mathcal{O}(g^2) ]}{\frac{1}{a} \exp(-\int \frac{dg}{\beta})} = 1 + \mathcal{O}(g^2).
\end{equation}
The exponential integral $\exp(-\int dg/\beta)$, which carries the essential singularity, is common to both quantities and divides out identically. The remaining quotient is an asymptotic series with non-zero radius of convergence in the effective coupling, hence real-analytic.
\end{proof}

\section{Component III: Positivity as the Sole Non-Analytic Content}
\label{sec:factor_positivity}

Because the existence factor $\mathcal{E} = 1$ is guaranteed by reflection positivity, and the value factor $\mathcal{V}$ is analytic, the entire non-analytic content of the Yang--Mills mass gap problem is concentrated in the third factor:
\begin{equation}
\mathcal{P} \in \{0, 1\}, \qquad \mathcal{P} = \begin{cases} 1 & \text{if } \Delta > 0, \\ 0 & \text{if } \Delta = 0. \end{cases}
\end{equation}

\begin{theorem}[Reduction of Positivity to the Centre-Charged Sector]
\label{thm:positivity_reduction}
The positivity factor $\mathcal{P} = 1$ is mathematically equivalent to the statement that the transfer matrix operator restricted to the centre-charged superselection sector $\mathcal{H}^{(1)}$ satisfies a uniform spectral gap:
\begin{equation}
\label{eq:centre_positivity_condition}
\mathcal{P} = 1 \iff \inf_{a > 0} \frac{1}{a} \left( 1 - \|\hat{T}^{(1)}_a\|_{\op} \right) > 0.
\end{equation}
\end{theorem}

The physical content of \eqref{eq:centre_positivity_condition} is the \textbf{persistence of color confinement}:
\begin{itemize}
    \item The sector $\mathcal{H}^{(1)}$ carries the chromoelectric flux of fundamental representation $R = \mathbf{N}$ under the centre $Z(\SU(N)) \cong \mathbb{Z}_N$.
    \item The energy of the lowest state in $\mathcal{H}^{(1)}$ on a spatial torus of size $L$ is $E_1(L) = \sigma \, L$, where $\sigma$ is the confining string tension.
    \item If the string tension vanishes ($\sigma = 0$), the centre symmetry is broken, the confining flux delocalizes, and $\mathcal{P} = 0$, producing a gapless deconfined phase (as occurs in the $\U(1)$ Coulomb phase).
    \item If the string tension is strictly positive ($\sigma > 0$), the centre symmetry remains unbroken, the confining flux forms a stable flux tube, and the lowest scalar glueball state in $\mathcal{H}^{(0)}$ satisfies:
    \begin{equation}
    M_{0^{++}} \ge 2 \sqrt{\sigma} > 0 \implies \mathcal{P} = 1.
    \end{equation}
\end{itemize}

\section{Meeting the Analyticity Obstruction}
\label{sec:meeting_obstruction}

The factorisation $\Delta = \mathcal{E} \cdot \mathcal{V} \cdot \mathcal{P}$ resolves the historical impasse of four-dimensional constructive field theory:
1. It explains why perturbation theory failed: perturbation theory can compute the analytic value factor $\mathcal{V}$, but is blind to the non-analytic positivity factor $\mathcal{P}$.
2. It meets the analyticity obstruction not by attempting to expand around the essential singularity, but by **isolating the non-analyticity into a single discrete symmetry sector**: the confining string tension of the centre-charged sector $\mathcal{H}^{(1)}$.
3. It reduces the mass gap problem to proving that the confining flux does not collapse across the marginal crossover window $\beta \in [1/2, 5/2]$.

This reduction converts the mass gap problem from an intractable infinite-dimensional analytic problem into a finite, geometric problem of controlling the discrete centre flux across a compact window.
